En este Trabajo de Fin de Máster se estudia el problema de los factores compartidos en claves públicas RSA extraídas de \textit{Certificate Transparency logs}. Esta vulnerabilidad puede aparecer cuando la generación de claves no se realiza con suficiente aleatoriedad, provocando que distintos módulos RSA compartan un mismo factor primo y puedan factorizarse de forma eficiente.
\\
\\
El trabajo se apoya, por un lado, en la tesis \textit{Finding shared RSA factors in the Certificate Transparency logs}, que abordó este problema mediante un enfoque clásico de \textit{batch GCD}, y, por otro, en el artículo de Pelofske, que propone el algoritmo \textit{Binary Tree Batch GCD} como una alternativa más eficiente. A partir de estas referencias, se desarrolla una implementación propia en C++17 con GMP, se compara frente a las implementaciones Python originales del artículo y se aplica sobre módulos RSA reales extraídos de \textit{CT logs}.
\\
\\
La evaluación experimental reproduce de forma reducida la estructura del experimento de Pelofske, utilizando módulos RSA de 1024 y 2048 bits, tres niveles de claves débiles y distintos tamaños de entrada. Los resultados confirman que el algoritmo \textit{Binary Tree Batch GCD} mejora de forma consistente al método clásico basado en \textit{remainder tree}, y muestran que la implementación C++/GMP reduce de manera significativa tanto el tiempo de ejecución como el consumo de memoria.
\\
\\
Además, el procedimiento se aplica a un conjunto real de módulos obtenidos a partir de \textit{ct-archive} y de los datos recopilados durante el trabajo. La extracción se centra en los certificados de emisores disponibles en los archivos ZIP de los logs, obtiene los módulos RSA con OpenSSL, normaliza los valores en hexadecimal y elimina duplicados antes de ejecutar el ataque. El conjunto final contiene 54.278 módulos únicos brutos, de los cuales 54.275 se consideran módulos RSA plausibles tras la validación. Sobre este conjunto no se detectan factores compartidos no triviales ni claves vulnerables, y la ejecución completa del algoritmo C++/GMP tarda aproximadamente 23,8 segundos en el entorno local utilizado. En conjunto, el TFM combina estudio de trabajos previos, desarrollo algorítmico, implementación experimental y validación aplicada sobre un problema real de ciberseguridad.
